% Generisano: uređuj beleske/sNNN.tex i naslove u manifestu.
\documentclass[11pt,a4paper]{article}
% Zajednički preambulum za dokument za čitanje; nije Beamer tema.
\usepackage[a4paper,margin=20mm,headheight=15pt]{geometry}
\usepackage{lmodern,fontspec}
\usepackage[serbian]{babel}
\setmainfont{Latin Modern Sans}
\setsansfont{Latin Modern Sans}
\setmonofont{Latin Modern Mono}
\usepackage{amsmath,amssymb,mathtools,graphicx,xcolor}
\usepackage{array,booktabs,tabularx,multirow,makecell,listings,fancyhdr}
\pagestyle{fancy}\fancyhf{}
\newcommand{\BeleskeTekuciID}{NEOZNACEN}
\fancyhead[L]{\small Beleške uz slajd \texttt{\expandafter\detokenize\expandafter{\BeleskeTekuciID}}}
\fancyfoot[R]{\small\thepage}
\setlength{\parindent}{0pt}\setlength{\parskip}{6pt}
\newwrite\BeleskeMapa
\immediate\openout\BeleskeMapa=\jobname.notes-map.tsv
\AddToHook{shipout/before}{%
  \immediate\write\BeleskeMapa{\BeleskeTekuciID\space\thepage}}
\AddToHook{enddocument/afterlastpage}{\immediate\closeout\BeleskeMapa}
\newcommand{\BeleskeID}[1]{%
  \def\BeleskeZadatiID{#1}%
  \ifx\BeleskeTekuciID\BeleskeZadatiID\else
    \PackageError{beleske}{Pogresan ID beleske: #1}{Proveri mapiranje slajda.}%
  \fi}
\newcommand{\BeleskaSlajda}[4]{%
  \clearpage\gdef\BeleskeTekuciID{#1}%
  {\Large\bfseries Slajd \textnormal{\texttt{\detokenize{#1}}}: #3\par}\medskip
  \includegraphics[page=#2,width=\linewidth]{\BeleskePrezentacija}\par\medskip
  \input{#4}}

\newcommand{\BeleskePrezentacija}{build/10_logicka_kola_sa_bipolarnim_tranzistorima.pdf}
\begin{document}
\BeleskaSlajda{10-s001}{1}{Logička kola sa bipolarnim tranzistorima}{beleske/s001.tex}
\BeleskaSlajda{10-s002}{2}{Model diode u logičkim kolima}{beleske/s002.tex}
\BeleskaSlajda{10-s003}{3}{Zener i Šotki dioda}{beleske/s003.tex}
\BeleskaSlajda{10-s004}{4}{Impulsni režim rada diode}{beleske/s004.tex}
\BeleskaSlajda{10-s005}{5}{Model bipolarnog tranzistora}{beleske/s005.tex}
\BeleskaSlajda{10-s006}{6}{Režimi bipolarnog tranzistora}{beleske/s006.tex}
\BeleskaSlajda{10-s007}{7}{Određivanje režima rada}{beleske/s007.tex}
\BeleskaSlajda{10-s008}{8}{Impulsni režim bipolarnog tranzistora}{beleske/s008.tex}
\BeleskaSlajda{10-s009}{9}{Kako ubrzati rad}{beleske/s009.tex}
\BeleskaSlajda{10-s010}{10}{RTL logička kola}{beleske/s010.tex}
\BeleskaSlajda{10-s011}{11}{Karakteristika prenosa: početak}{beleske/s011.tex}
\BeleskaSlajda{10-s012}{12}{Ivica provođenja}{beleske/s012.tex}
\BeleskaSlajda{10-s013}{13}{Ulazak u aktivni režim}{beleske/s013.tex}
\BeleskaSlajda{10-s014}{14}{Aktivna karakteristika RTL kola}{beleske/s014.tex}
\BeleskaSlajda{10-s015}{15}{Diskretni model i prekid karakteristike}{beleske/s015.tex}
\BeleskaSlajda{10-s016}{16}{Nagib aktivne oblasti}{beleske/s016.tex}
\BeleskaSlajda{10-s017}{17}{Zasićenje i margine šuma}{beleske/s017.tex}
\BeleskaSlajda{10-s018}{18}{Prag prebacivanja i margine}{beleske/s018.tex}
\BeleskaSlajda{10-s019}{19}{Drugi način određivanja visokog praga}{beleske/s019.tex}
\BeleskaSlajda{10-s020}{20}{Razlika diskretnih pragova}{beleske/s020.tex}
\BeleskaSlajda{10-s021}{21}{Ulazni strujni kapacitet RTL kola}{beleske/s021.tex}
\BeleskaSlajda{10-s022}{22}{Kompromis pri izboru baznog otpornika}{beleske/s022.tex}
\BeleskaSlajda{10-s023}{23}{Izlazni strujni kapacitet jedinice}{beleske/s023.tex}
\BeleskaSlajda{10-s024}{24}{Maksimalna struja i izbor kolektorskog otpornika}{beleske/s024.tex}
\BeleskaSlajda{10-s025}{25}{Izlazni strujni kapacitet nule}{beleske/s025.tex}
\BeleskaSlajda{10-s026}{26}{Nulti strujni kapacitet na granici}{beleske/s026.tex}
\BeleskaSlajda{10-s027}{27}{Faktor grananja RTL kola}{beleske/s027.tex}
\BeleskaSlajda{10-s028}{28}{Povećanje strujnog kapaciteta}{beleske/s028.tex}
\BeleskaSlajda{10-s029}{29}{Dinamičke karakteristike: punjenje}{beleske/s029.tex}
\BeleskaSlajda{10-s030}{30}{Dinamičke karakteristike: pražnjenje}{beleske/s030.tex}
\BeleskaSlajda{10-s031}{31}{Trenutak ulaska u zasićenje}{beleske/s031.tex}
\BeleskaSlajda{10-s032}{32}{Kašnjenje i vreme opadanja RTL izlaza}{beleske/s032.tex}
\BeleskaSlajda{10-s033}{33}{Integrisano RTL kolo}{beleske/s033.tex}
\BeleskaSlajda{10-s034}{34}{Višeulazna RTL kola}{beleske/s034.tex}
\BeleskaSlajda{10-s035}{35}{Logička kola sa otvorenim kolektorom}{beleske/s035.tex}
\BeleskaSlajda{10-s036}{36}{Otvoreni kolektor bez podiznog otpornika}{beleske/s036.tex}
\BeleskaSlajda{10-s037}{37}{Prikaz signala i visoke impedanse}{beleske/s037.tex}
\BeleskaSlajda{10-s038}{38}{Neiskorišćeni ulazi i terminacija}{beleske/s038.tex}
\BeleskaSlajda{10-s039}{39}{Diodna ILI funkcija: oba ulaza niska}{beleske/s039.tex}
\BeleskaSlajda{10-s040}{40}{Diodna ILI funkcija: jedan visok ulaz}{beleske/s040.tex}
\BeleskaSlajda{10-s041}{41}{Diodna ILI funkcija: zaključak}{beleske/s041.tex}
\BeleskaSlajda{10-s042}{42}{Diodna I funkcija}{beleske/s042.tex}
\BeleskaSlajda{10-s043}{43}{DTL logička kola}{beleske/s043.tex}
\BeleskaSlajda{10-s044}{44}{Uklanjanje negativnog napajanja}{beleske/s044.tex}
\BeleskaSlajda{10-s045}{45}{DTL: analiza na invertoru}{beleske/s045.tex}
\BeleskaSlajda{10-s046}{46}{Karakteristika prenosa DTL kola}{beleske/s046.tex}
\BeleskaSlajda{10-s047}{47}{DTL: pojačanje i margine šuma}{beleske/s047.tex}
\BeleskaSlajda{10-s048}{48}{Strujni kapacitet DTL kola}{beleske/s048.tex}
\BeleskaSlajda{10-s049}{49}{Dinamičke karakteristike DTL kola}{beleske/s049.tex}
\BeleskaSlajda{10-s050}{50}{TTL logička kola}{beleske/s050.tex}
\BeleskaSlajda{10-s051}{51}{Pomoćni prikaz TTL ulaznog tranzistora}{beleske/s051.tex}
\BeleskaSlajda{10-s052}{52}{Inverzni režim rada tranzistora}{beleske/s052.tex}
\BeleskaSlajda{10-s053}{53}{Višeulazna TTL kola}{beleske/s053.tex}
\BeleskaSlajda{10-s054}{54}{Prelaz na standardno TTL kolo}{beleske/s054.tex}
\BeleskaSlajda{10-s055}{55}{Standardno TTL logičko kolo}{beleske/s055.tex}
\BeleskaSlajda{10-s056}{56}{Zaštita ulaza i totem-pole izlaz}{beleske/s056.tex}
\BeleskaSlajda{10-s057}{57}{Karakteristika prenosa TTL kola}{beleske/s057.tex}
\BeleskaSlajda{10-s058}{58}{Neopterećeni totem-pole izlaz}{beleske/s058.tex}
\BeleskaSlajda{10-s059}{59}{TTL: naponski nivoi i margine}{beleske/s059.tex}
\BeleskaSlajda{10-s060}{60}{Strujni kapacitet TTL ulaza}{beleske/s060.tex}
\BeleskaSlajda{10-s061}{61}{Strujni kapacitet TTL izlaza}{beleske/s061.tex}
\BeleskaSlajda{10-s062}{62}{Totem pole: granica aktivnog režima}{beleske/s062.tex}
\BeleskaSlajda{10-s063}{63}{Punjenje izlaza: aktivni režim}{beleske/s063.tex}
\BeleskaSlajda{10-s064}{64}{Punjenje izlaza: zasićeni režim}{beleske/s064.tex}
\BeleskaSlajda{10-s065}{65}{Trenutak izlaska iz zasićenja}{beleske/s065.tex}
\BeleskaSlajda{10-s066}{66}{Dinamički režim TTL kola}{beleske/s066.tex}
\BeleskaSlajda{10-s067}{67}{Low power TTL}{beleske/s067.tex}
\BeleskaSlajda{10-s068}{68}{High power TTL}{beleske/s068.tex}
\BeleskaSlajda{10-s069}{69}{Schottky TTL}{beleske/s069.tex}
\BeleskaSlajda{10-s070}{70}{Schottky tranzistor}{beleske/s070.tex}
\BeleskaSlajda{10-s071}{71}{Schottky TTL sa aktivnim pražnjenjem}{beleske/s071.tex}
\BeleskaSlajda{10-s072}{72}{Low power Schottky TTL}{beleske/s072.tex}
\BeleskaSlajda{10-s073}{73}{Kataloški dozvoljeni opsezi}{beleske/s073.tex}
\BeleskaSlajda{10-s074}{74}{Kataloške DC i AC karakteristike}{beleske/s074.tex}
\BeleskaSlajda{10-s075}{75}{Pakovanje i oznake: 74LS00}{beleske/s075.tex}
\BeleskaSlajda{10-s076}{76}{Spajanje totem-pole izlaza}{beleske/s076.tex}
\BeleskaSlajda{10-s077}{77}{Zajednička magistrala i dozvola rada}{beleske/s077.tex}
\BeleskaSlajda{10-s078}{78}{Trostatički TTL izlaz}{beleske/s078.tex}
\BeleskaSlajda{10-s079}{79}{Kašnjenja prelaza u visoku impedansu}{beleske/s079.tex}
\BeleskaSlajda{10-s080}{80}{LS125A i LS126A: šema i kašnjenja}{beleske/s080.tex}
\BeleskaSlajda{10-s081}{81}{ECL — Emitter Coupled Logic}{beleske/s081.tex}
\BeleskaSlajda{10-s082}{82}{ECL: diferencijalni par}{beleske/s082.tex}
\BeleskaSlajda{10-s083}{83}{ECL: mali ulazni napon}{beleske/s083.tex}
\BeleskaSlajda{10-s084}{84}{ECL: preusmeravanje struje}{beleske/s084.tex}
\BeleskaSlajda{10-s085}{85}{ECL: granica zasićenja}{beleske/s085.tex}
\BeleskaSlajda{10-s086}{86}{ECL: prenosna karakteristika}{beleske/s086.tex}
\BeleskaSlajda{10-s087}{87}{ECL: izbor referentnog napona}{beleske/s087.tex}
\BeleskaSlajda{10-s088}{88}{ECL: precizan izbor elemenata}{beleske/s088.tex}
\BeleskaSlajda{10-s089}{89}{ECL: izlazni baferi}{beleske/s089.tex}
\BeleskaSlajda{10-s090}{90}{ECL: napajanje i smetnje}{beleske/s090.tex}
\BeleskaSlajda{10-s091}{91}{ECL: strujni kapaciteti}{beleske/s091.tex}
\BeleskaSlajda{10-s092}{92}{ECL: dinamički režim}{beleske/s092.tex}
\BeleskaSlajda{10-s093}{93}{Standardno dvoulazno ECL kolo}{beleske/s093.tex}
\BeleskaSlajda{10-s094}{94}{ECL: PECL DC karakteristike}{beleske/s094.tex}
\BeleskaSlajda{10-s095}{95}{ECL: NECL DC karakteristike}{beleske/s095.tex}
\BeleskaSlajda{10-s096}{96}{ECL: AC karakteristike}{beleske/s096.tex}
\end{document}
